\documentclass[../../../include/open-logic-section]{subfiles}

\begin{document}

\olfileid{sth}{choice}{hartogs}
\olsection{Kabandingan lan Lema Hartogs}

Mau sisih kang nguntungaké. Saiki sisih kang nyusahaké. Tanpa
Aksioma Pilihan, prakarané dadi \emph{ruwet}. Kanggo ndeleng sebabé,
gatekna asil saka \cite{Hartogs1915} iki:

\begin{lem}[\emph{ing $\ZF$}]\ollabel{HartogsLemma}
Kanggo sembarang himpunan $A$, ana ordinal $\alpha$ kang nyukupi
$\cardnless{\alpha}{A}$.
\end{lem}

\begin{proof}
Yèn $B \subseteq A$ lan $R \subseteq B^2$, mula $\tuple{B, R} \subseteq
V_{\setrank{A}+4}$ miturut
\olref[ord-arithmetic][using-addition]{rankcomputation}. Mula, kanthi
Pamisahan, gatekna:
\[
	C = \Setabs{\tuple{B, R} \in V_{\setrank{A}+5}}{B\subseteq A
	\text{ lan $\tuple{B, R}$ sawijining pengurutan becik}}
\]
Kanthi Panggantian lan
\olref[ordinals][ordtype]{thmOrdinalRepresentation}, wujudna himpunan:
\[
	\alpha = \Setabs{\ordtype{B, R}}{\tuple{B, R} \in C}.
\]
Miturut \olref[ordinals][basic]{corordtransitiveord}, $\alpha$ iku
ordinal amarga sawijining himpunan transitif saka ordinal. Pancen, yèn
$\gamma \in \beta \in \alpha$, mula $\beta = \ordtype{B, R}$ kanggo
sawijining $B \subseteq A$; saka kéné $\gamma = \ordtype{B_b, R_b}$
kanggo sawijining $b \in B$ miturut
\olref[ordinals][iso]{wellordinitialsegment}, saéngga
$\gamma \in \alpha$.

Kanggo bukti kanthi kontradiksi, anggepen ana !!a{injection}
$f \colon \alpha \to A$. Banjur tetepna:
\begin{align*}
	B &= \ran{f}\\
	R &= \Setabs{\tuple{f(\xi), f(\zeta)}}{\xi,\zeta\in\alpha \land \xi\in\zeta}.
\end{align*}
Cetha yèn $\alpha = \ordtype{B, R}$ lan $\tuple{B, R} \in C$.
Mula $\alpha \in \alpha$, sawijining kontradiksi.
\end{proof}

Asil iki mènèhi téoréma kang jero:

\begin{thm}[\emph{ing $\ZF$}]
Pratelan-pratelan iki ekuivalen:
\begin{enumerate}
	\item\ollabel{equivwo} Aksioma Pengurutan Becik
	\item\ollabel{equivcompare} $\cardle{A}{B}$ utawa
	$\cardle{B}{A}$ kanggo sembarang himpunan $A$ lan $B$
\end{enumerate}
\end{thm}

\begin{proof}
\emph{\olref{equivwo} $\Rightarrow$ \olref{equivcompare}.} Tetepna $A$ lan
$B$. Miturut \olref{equivwo}, ana pengurutan becik $\tuple{A, R}$
lan $\tuple{B, S}$. Miturut
\olref[ordinals][ordtype]{thmOrdinalRepresentation}, jupuken
isomorfisme $f \colon \alpha \to A$ lan $g \colon \beta \to B$.
Miturut \olref[sth][ordinals][basic]{ordinalsaresubsets}, salah siji
$\alpha \subseteq \beta$ utawa $\beta \subseteq \alpha$.
Yèn $\alpha \subseteq \beta$, mula
$\comp{f^{-1}}{g} \colon A \to B$ iku !!a{injection}, saéngga
$\cardle{A}{B}$; semono uga, yèn $\beta \subseteq \alpha$, mula
$\cardle{B}{A}$.

\emph{\olref{equivcompare} $\Rightarrow$ \olref{equivwo}.} Tetepna $A$;
miturut \olref{HartogsLemma} ana ordinal $\beta$ kang nyukupi
$\cardnless{\beta}{A}$. Kanthi \olref{equivcompare}, kita nduwèni
$\cardle{A}{\beta}$. Mula ana !!{injection} $f \colon A \to
\beta$, lan injeksi iki bisa digunakaké kanggo ngurutaké unsur
$A$ kanthi becik, kanthi ndhéfinisikaké pengurutan
$\Setabs{\tuple{a, b} \in A \times A}{f(a) \in f(b)}$.
\end{proof}
\noindent
Akibat langsungé: yèn Pengurutan Becik gagal, ana himpunan-himpunan
kang \emph{pancen ora bisa dibandhingaké} ukurané. Mula aritmetika
kardinal transfinit bakal ruwet. Upamané, kita kudu ninggalaké
pratelan umum yèn kanggo $A$ lan $B$ tanpa wates,
$\cardeq{A \disjointsum B}{M}$ lan $\cardeq{A \times B}{M}$,
ing ngendi $M$ iku kang luwih gedhé saka $A$ lan $B$ (delengen
\olref[card-arithmetic][simp]{cardplustimesmax}). Masalahé prasaja:
yèn ukuran $A$ lan $B$ ora bisa \emph{dibandhingaké}, takon endi
kang luwih gedhé dadi ora ana tegesé.

\end{document}
